\documentclass[11pt]{article}

%-------------------------
% Packages
%-------------------------
\usepackage[a4paper,margin=1in]{geometry}
\usepackage{xcolor}
\usepackage{tikz}
\usepackage[most]{tcolorbox}
\usepackage{mathtools}
\usepackage{amsthm, amsmath, amssymb}
\usepackage{enumitem}
\usepackage{hyperref}
\usepackage[nameinlink,noabbrev]{cleveref}
\usepackage{titling} 
\usepackage{fontspec}
\usepackage{unicode-math}


%-------------------------
% Fonts
%-------------------------
\setmainfont{JetBrains Mono}
\setsansfont{JetBrains Mono}
\setmonofont{JetBrains Mono}


\setmathfont[
    Scale=1.25
]{latinmodern-math.otf}

\setmathfont[
  range=\mathup/{latin, Latin, 0-9},
  range=\mathit/{latin, Latin},
  range=\mathbfit/{latin, Latin},
  Scale=1.5
]{JetBrains Mono Bold} 

\setmathfont[range=\mathbb, Scale=MatchUppercase]{latinmodern-math.otf}

%-------------------------
% Paragraph spacing
%-------------------------
\setlength{\parindent}{0pt}
\setlength{\parskip}{0.7em}


%-------------------------
% Dark mode colors
%-------------------------
\definecolor{bgcolor}{HTML}{333333}
\definecolor{textcolor}{HTML}{FAFAFA}
\definecolor{defcolor}{HTML}{E86873}
\definecolor{thmcolor}{HTML}{0A9396}
\definecolor{lemcolor}{HTML}{94D2BD}
\definecolor{corcolor}{HTML}{9B4AF7}
\definecolor{probcolor}{HTML}{EE9B00}
\definecolor{excolor}{HTML}{21E933}
\definecolor{propcolor}{HTML}{EE9B20}

%-------------------------
% Background and text color
%-------------------------
\pagecolor{bgcolor}
\color{textcolor}

%-------------------------
% Paragraph spacing
%-------------------------
\setlength{\parindent}{0pt}
\setlength{\parskip}{0.7em}

%-------------------------
% Base tcolorbox settings (Removed 'attach title to upper')
%-------------------------
\tcbset{
  enhanced,
  colback=bgcolor,
  colframe=thmcolor,
  coltext=white,
  coltitle=white,
  fonttitle=\bfseries,
  boxrule=0.7pt,
  left=1em,
  right=1em,
  top=0.7em,
  bottom=0.7em,
  before skip=10pt,
  after skip=10pt,
  title after break,
  breakable
}

%-------------------------
% Standard Theorem Environments (using amsthm)
% Syntax: \begin{thm}[Optional Title]\label{thm:label} ...
%-------------------------
\newtheorem{thm}{Theorem}[section]
\newtheorem*{thm*}{Theorem}
\newtheorem{defn}{Definition}[section]
\newtheorem*{defn*}{Definition}
\newtheorem{lem}{Lemma}[section]
\newtheorem*{lem*}{Lemma}
\newtheorem{cor}{Corollary}[section]
\newtheorem*{cor*}{Corollary}
\newtheorem{prob}{Problem}[section]
\newtheorem*{prob*}{Problem}
\newtheorem{ex}{Example}[section]
\newtheorem*{ex*}{Example}
\newtheorem{prop}{Proposition}[section]
\newtheorem*{prop*}{Proposition}

%-------------------------
% Tcolorbox environments (applies style to amsthm environments)
% Removed 'attach title to upper' from all individual settings.
%-------------------------
\tcolorboxenvironment{thm}{
  colframe=thmcolor,
  colback=thmcolor!15!bgcolor,
  fonttitle=\bfseries,
  coltitle=white,
  breakable,
  title after break,
}

\tcolorboxenvironment{thm*}{
  colframe=thmcolor,
  colback=thmcolor!15!bgcolor,
  fonttitle=\bfseries,
  coltitle=white,
  breakable,
  title after break,
}

\tcolorboxenvironment{defn}{
  colframe=defcolor,
  colback=defcolor!15!bgcolor,
  fonttitle=\bfseries,
  coltitle=white,
  breakable,
  title after break,
}

\tcolorboxenvironment{defn*}{
  colframe=defcolor,
  colback=defcolor!15!bgcolor,
  fonttitle=\bfseries,
  coltitle=white,
  breakable,
  title after break,
}

\tcolorboxenvironment{lem}{
  colframe=lemcolor,
  colback=lemcolor!15!bgcolor,
  fonttitle=\bfseries,
  coltitle=white,
  breakable,
  title after break,
}

\tcolorboxenvironment{lem*}{
  colframe=lemcolor,
  colback=lemcolor!15!bgcolor,
  fonttitle=\bfseries,
  coltitle=white,
  breakable,
  title after break,
}

\tcolorboxenvironment{cor}{
  colframe=corcolor,
  colback=corcolor!15!bgcolor,
  fonttitle=\bfseries,
  coltitle=white,
  breakable,
  title after break,
}

\tcolorboxenvironment{cor*}{
  colframe=corcolor,
  colback=corcolor!15!bgcolor,
  fonttitle=\bfseries,
  coltitle=white,
  breakable,
  title after break,
}

\tcolorboxenvironment{prob}{
  colframe=probcolor,
  colback=probcolor!15!bgcolor,
  fonttitle=\bfseries,
  coltitle=white,
  breakable,
  title after break,
}

\tcolorboxenvironment{prob*}{
  colframe=probcolor,
  colback=probcolor!15!bgcolor,
  fonttitle=\bfseries,
  coltitle=white,
  breakable,
  title after break,
}

\tcolorboxenvironment{ex}{
  colframe=excolor,
  colback=excolor!15!bgcolor,
  fonttitle=\bfseries,
  coltitle=white,
  breakable,
  title after break,
}

\tcolorboxenvironment{ex*}{
  colframe=excolor,
  colback=excolor!15!bgcolor,
  fonttitle=\bfseries,
  coltitle=white,
  breakable,
  title after break,
}

\tcolorboxenvironment{prop}{
    colframe=propcolor, 
    colback=propcolor!15!bgcolor,
    fonttitle=\bfseries,
    coltitle=white,
    breakable,
    title after break
}

\tcolorboxenvironment{prop*}{
    colframe=propcolor, 
    colback=propcolor!15!bgcolor,
    fonttitle=\bfseries,
    coltitle=white,
    breakable,
    title after break
}

%-------------------------
% Cleveref settings (Using short environment names for correct display)
%-------------------------
\crefname{thm}{theorem}{theorems}
\Crefname{thm}{Theorem}{Theorems}

\crefname{defn}{definition}{definitions}
\Crefname{defn}{Definition}{Definitions}

\crefname{lem}{lemma}{lemmas}
\Crefname{lem}{Lemma}{Lemmas}

\crefname{cor}{corollary}{corollaries}
\Crefname{cor}{Corollary}{Corollaries}

\crefname{prob}{problem}{problems}
\Crefname{prob}{Problem}{Problems}

\crefname{ex}{example}{examples}
\Crefname{ex}{Example}{Examples}

\crefname{prop}{proposition}{propositions}
\Crefname{prop}{Proposition}{Propositions}

\crefname{thm*}{theorem}{theorems}
\Crefname{thm*}{Theorem}{Theorems}

\crefname{defn*}{definition}{definitions}
\Crefname{defn*}{Definition}{Definitions}

\crefname{lem*}{lemma}{lemmas}
\Crefname{lem*}{Lemma}{Lemmas}

\crefname{cor*}{corollary}{corollaries}
\Crefname{cor*}{Corollary}{Corollaries}

\crefname{prob*}{problem}{problems}
\Crefname{prob*}{Problem}{Problems}

\crefname{ex*}{example}{examples}
\Crefname{ex*}{Example}{Examples}

\crefname{prop*}{proposition}{propositions}
\Crefname{prop*}{Proposition}{Propositions}


\title{\huge{Extra study and consolidation in Algebra}}
\author{\LARGE{Thobias Høivik}}
\date{}

\begin{document}
\maketitle

\newpage
\tableofcontents

\newpage

\section{What is this?}
Due to my being very interested in algebra, and interested 
in becoming very good at it, I have these extra 
notes for me to work through a lot of 
problems. 
Basically, there will probably not be any pedagogical value here. 

\section{Repeating some stuff about ideals}

\begin{prob}
    Let 
    \[
        R = \frac{k[x,y]}{(xy)}
    \]
    and write $\overline x, \overline y$ for the residue classes. 

    What is $(\overline x)$ as an ideal of $R$? 

    Compute 
    \[
        R/(\bar x),\qquad R/(\bar y),\qquad R/(\bar x,\bar y).
    \]
    Then classify them as prime/maximal/neither. 
\end{prob}
\begin{proof}[Sol]
    We note that the ideal 
    generated by $\overline x = x + (xy)$ is 
    \[
        (\overline x) = 
        \{\overline{xf} \bigm| f \in k[x,y]\}. 
    \]
    We have $(xy) \subseteq (x)$ and 
    $(xy) \subseteq (y)$ as ideals of $k[x,y]$. 

    Consider the projection $\pi : k[x,y] \to R$ 
    by 
    \[
        \pi(a) = \overline a = a + (xy). 
    \]
    Since $\pi$ is surjective and $\{x\} \subseteq k[x,y]$ we have 
    \[
        \pi[(x)] = (\pi x) = (\overline x),  
    \]
    (note: surjectivity guarantees that scalar multiplication in 
    $R$ can be pulled back to $k[x,y]$, ensuring 
    $\pi[(x)] \supseteq (\pi(x))$)
    but also 
    \[
        \pi[(x)] = \frac{(x)}{(xy)}
    \]
    since $(xy) \subseteq (x)$. 

    Similarly,
    \[
        \frac{(y)}{(xy)} = \pi[(y)] = (\overline y),
    \]
    and 
    \[
        \frac{(x,y)}{(xy)} = \pi[(x,y)] = (\overline x, \overline y).
    \]

    Thus we can apply the third isomorphism theorem to get 
    \[
        \frac{R}{(\overline x)} = \frac{k[x,y]/(xy)}{(x)/(xy)} 
        \cong \frac{k[x,y]}{(x)} \cong k[y], 
    \]
    and 
    \[
        \frac{R}{(\overline y)} = \frac{k[x,y]/(xy)}{(y)/(xy)} 
        \cong \frac{k[x,y]}{(y)} \cong k[x], 
    \]
    and 
    \[
        \frac{R}{(\overline x, \overline y)} = 
        \frac{k[x,y]/(xy)}{(x, y)/(xy)} 
        \cong \frac{k[x,y]}{(x, y)} \cong k.  
    \]

    The first two are integral domains so 
    $(\overline x), (\overline y)$ are prime, though not 
    maximal, while 
    $(\overline x, \overline y)$ is maximal since 
    $R/(\overline x, \overline y) \cong k$ is a field. 
\end{proof}

\begin{prob}
    Let 
    \[
        R = \frac{k[x,y]}{(xy)}, 
    \]
    and let 
    \[
        \pi : k[x,y] \to R
    \]
    be the quotient map. 

    Prove that every prime ideal $P \subseteq R$ 
    contains $\overline x$ or $\overline y$. 

    Then, classify every prime ideal of $R$. 

    You may use without proof: 
    \[
        P \subseteq R \ \text{prime} 
        \Leftrightarrow \pi^{-1}(P) \subseteq k[x,y] 
        \ \text{is prime containing} \ (xy). 
    \]
\end{prob}

\begin{proof}
    We have $\overline x \overline y = 0$, and since 
    $P$ is an ideal, and thus a group, we have 
    $0 \in P$, and since $P$ is prime we must then have 
    \[\overline x \in P \ \text{or} \ \overline y \in P.\]

    Suppose $P$ contains $\overline x$. 
    Then, $(\overline x) \subseteq P$. 

    Then, by the third isomorphism theorem,  
    \[\frac{R/(\overline x)}{P / (\overline x)}         \cong R/P,\]
    but 
    \[\frac{R}{(\overline x)} \cong k[y]\]
    so 
    \[R/P \cong \frac{k[y]}{P/(\overline x)}.\]

    Since $P$ is a prime ideal of $R$, the quotient $R/P$ is an integral domain. Therefore, $P/(\overline x)$ must be a prime ideal of $k[y]$. 

    Since $k[y]$ is a Principal Ideal Domain (PID), its prime ideals are:
    \begin{enumerate}
        \item The zero ideal $(0)$, or
        \item Irreducible polynomial ideals $(f(y))$, where $f(y) \in k[y]$ is an irreducible polynomial.
    \end{enumerate}

    Under the canonical quotient map $R \to R/(\overline x) \cong k[y]$, these correspond precisely to the prime ideals of $R$ containing $\overline x$:
    \begin{enumerate}
        \item $P = (\overline x)$, corresponding to $(0) \subset k[y]$.
        \item $P = (\overline x, f(\overline y))$, where $f(y) \in k[y]$ is irreducible.
    \end{enumerate}

    By symmetric reasoning, if $\overline y \in P$, then $P/(\overline y)$ is a prime ideal of $k[x]$, which yields:
    \begin{enumerate}
        \item $P = (\overline y)$, corresponding to $(0) \subset k[x]$.
        \item $P = (\overline y, g(\overline x))$, where $g(x) \in k[x]$ is irreducible.
    \end{enumerate}

    Combining both cases, the prime ideals of $R$ consist of:
    \begin{itemize}
        \item Minimal prime ideals: $(\overline x)$ and $(\overline y)$.
        \item Maximal ideals: $(\overline x, f(\overline y))$ for irreducible $f(y) \in k[y]$, and $(\overline y, g(\overline x))$ for irreducible $g(x) \in k[x]$.
    \end{itemize}
\end{proof}

\end{document}
